%%=============================================================================
%% Conclusie
%%=============================================================================

\chapter{Conclusie}%
\label{ch:conclusie}

% TODO: Trek een duidelijke conclusie, in de vorm van een antwoord op de
% onderzoeksvra(a)g(en). Wat was jouw bijdrage aan het onderzoeksdomein en
% hoe biedt dit meerwaarde aan het vakgebied/doelgroep? 
% Reflecteer kritisch over het resultaat. In Engelse teksten wordt deze sectie
% ``Discussion'' genoemd. Had je deze uitkomst verwacht? Zijn er zaken die nog
% niet duidelijk zijn?
% Heeft het onderzoek geleid tot nieuwe vragen die uitnodigen tot verder 
%onderzoek?

Het beoogde resultaat is een onderbouwde en herhaalbare aanpak voor branching en deployment die geschikt is voor gebruik binnen de IDz-omgeving van Securex. Deze benadering dient heldere richtlijnen te geven voor het werken met sectoren, het aanbrengen van aanpassingen en het gecontroleerd toepassen van code in diverse omgevingen. Er wordt verwacht dat dit bijdraagt aan een verhoogde traceerbaarheid van de code, waardoor op elk moment inzicht kan verkregen worden in welke versie van de code in welke omgeving actief is.

